\SourceLeaf{393}%
geraden Linie weg; allein in diesem Fall ist leicht zu zeigen, dass nothwendig dieser gemeinschaftliche Halbirungspunkt zugleich der Mittelpunkt der Ellipse selbst sein wird.

2) Verlängert man zwei einander gegenüber liegende Seiten des Vierecks bis zu ihrem Durchschnitt und eben so die beiden andern, so darf man auch die zwischen diesen beiden Durchschnitts-Punkten enthaltene gerade Linie, als eine verschwindende die vier Seiten des Vierecks berührende Ellipse ansehen. Der Halbirungspunkt derselben muss also in eben der geraden Linie liegen, welche die Halbirungspunkte der beiden Diagonalen verbindet. Diese allgemeine Eigenschaft eines jeden Vierecks ist meines Wissens bisher noch nicht bemerkt; ich werde davon unten einen einfachen directen Beweis geben.

Um die Rechnungen noch mehr abzukürzen, will ich jetzt annehmen, dass man diese gerade Linie selbst zur Abscissen-Linie gewählt habe, und folglich \(\varphi=0\) sei. Der Anfangspunkt der Abscissen bleibt wie vorher willkürlich. Eben diese Bestimmung \(\varphi=0\) macht nun eine der vier Fundamental-Gleichungen entbehrlich, und wir haben also zur Bestimmung der vier unbekannten Grössen \(t,u,r,\psi\) theils die drei Gleichungen
\[
\begin{aligned}
2aa+t-4ar\cos A+rr\cos2A-u\cos2(A-\psi)&=0\\
2a'a'+t-4a'r\cos A'+rr\cos2A'-u\cos2(A'-\psi)&=0\\
2a''a''+t-4a''r\cos A''+rr\cos2A''-u\cos2(A''-\psi)&=0
\end{aligned}
\]
theils die Bedingung, dass der Inhalt der Ellipse, welchem offenbar das Product \(\alpha\beta\) proportional ist. und folglich auch \(4\alpha\alpha\beta\beta\) oder \((rr-t)^2-uu\) ein Maximum sein soll.

Setzt man Kürze halber \(rr-t=\theta\) und
\[
\begin{aligned}
b&=2(a-r\cos A)^2\\
b'&=2(a'-r\cos A')^2\\
b''&=2(a''-r\cos A'')^2
\end{aligned}
\]
so werden obige Gleichungen
\[
\begin{aligned}
\theta+u\cos2(A-\psi)&=b\\
\theta+u\cos2(A'-\psi)&=b'\\
\theta+u\cos2(A''-\psi)&=b''
\end{aligned}
\]
